% GENERATED FILE. Edit canonical lessons, then run npm run generate:books.
% canonical-source-hash: fnv1a64:91e7de136eae4a26
% canonical-lessons: MR-C438-rimjhim, MR-C438-musaldhar, MR-C438-valiv, MR-C438-barasne, MR-C438-bhijne

\chapter{Showers and Drizzle}
\label{ch:mr-showers-and-drizzle}
\hlchaptermodality{438}

\begin{chapteropening}
\textbf{By the end of this chapter:} \emph{I can talk about drizzle, downpours and the first storms of summer, and getting wet.}

Complete the last lesson of chapter 438: I can talk about drizzle, downpours and the first storms of summer, and getting wet.
\end{chapteropening}

\section[\texorpdfstring{\mr{रिमझिम} (rimjhim)}{रिमझिम (rimjhim)}]{\mr{रिमझिम} (rimjhim) — a drizzle}

\label{lesson:MR-C438-rimjhim}



\begin{quote}
Before the new one: say the Marathi for to freeze, then the Marathi for to melt.
\end{quote}

\subsection*{You'll want to know: \mr{रिमझिम}}
\textbf{\mr{रिमझिम}} — \emph{rimjhim} — \textquotedblleft{}a drizzle\textquotedblright{}.

\begin{quote}
\mr{सकाळपासून} \mr{रिमझिम} \mr{चालू} \mr{आहे}.
\end{quote}

\begin{quote}
\emph{sakāḷpāsūn rimjhim cālū āhe.}
\end{quote}

\begin{quote}
It has been drizzling since morning.
\end{quote}

Feminine.

Also an adverb: \mr{रिमझिम} \mr{पाऊस}, gentle rain.

\begin{cousinweb}
\textbf{\mr{रिमझिम}} (\emph{rimjhim}) is an echo word that imitates light, steady rain; its origin is not settled.
\end{cousinweb}

\begin{grammarlens}[title={What you've built}]
One more word for getting around.
\end{grammarlens}

\subsection*{Your turn}
\begin{itemize}
\raggedright
  \item \emph{Say it:} \emph{rimjhim}
  \item \emph{Say it:} \emph{rimjhim}, once more
  \item \emph{Say it:} \emph{sakāḷpāsūn rimjhim cālū āhe.} — It has been drizzling since morning.
  \item \emph{Recall:} say the Marathi for to freeze, then the Marathi for to melt, then say \emph{rimjhim} again
\end{itemize}

\begin{tcolorbox}[breakable,colback=teal!4,colframe=teal!35!black,title={Before you move on}]
What does \mr{रिमझिम} mean? (\textquotedblleft{}A drizzle\textquotedblright{}.) Say it once more.
\end{tcolorbox}
\section[\texorpdfstring{\mr{मुसळधार} (musaḷdhār)}{मुसळधार (musaḷdhār)}]{\mr{मुसळधार} (musaḷdhār) — torrential}

\label{lesson:MR-C438-musaldhar}



\begin{quote}
Before the new one: say the Marathi for to melt, then the Marathi for a drizzle.
\end{quote}

\subsection*{You'll want to know: \mr{मुसळधार}}
\textbf{\mr{मुसळधार}} — \emph{musaḷdhār} — \textquotedblleft{}torrential\textquotedblright{}.

\begin{quote}
\mr{काल} \mr{रात्री} \mr{मुसळधार} \mr{पाऊस} \mr{पडला}.
\end{quote}

\begin{quote}
\emph{kāl rātrī musaḷdhār pāūs paḍlā.}
\end{quote}

\begin{quote}
Last night it rained in torrents.
\end{quote}

\mr{मुसळ} is a pestle and \mr{धार} a stream: rain as thick as pestles.

\begin{cousinweb}
\textbf{\mr{मुसळधार}} (\emph{musaḷdhār}) is built from \textbf{\mr{मुसळ}}, \textquotedblleft{}a pestle\textquotedblright{}, inherited from Sanskrit \emph{musala} (CDIAL 10223), and \textbf{\mr{धार}}, \textquotedblleft{}a stream\textquotedblright{}.
\end{cousinweb}

\begin{grammarlens}[title={What you've built}]
One more word for getting around.
\end{grammarlens}

\subsection*{Your turn}
\begin{itemize}
\raggedright
  \item \emph{Say it:} \emph{musaḷdhār}
  \item \emph{Say it:} \emph{musaḷdhār}, once more
  \item \emph{Say it:} \emph{kāl rātrī musaḷdhār pāūs paḍlā.} — Last night it rained in torrents.
  \item \emph{Recall:} say the Marathi for to melt, then the Marathi for a drizzle, then say \emph{musaḷdhār} again
\end{itemize}

\begin{tcolorbox}[breakable,colback=teal!4,colframe=teal!35!black,title={Before you move on}]
What does \mr{मुसळधार} mean? (\textquotedblleft{}Torrential\textquotedblright{}.) Say it once more.
\end{tcolorbox}
\section[\texorpdfstring{\mr{वळीव} (vaḷīv)}{वळीव (vaḷīv)}]{\mr{वळीव} (vaḷīv) — a pre-monsoon storm}

\label{lesson:MR-C438-valiv}



\begin{quote}
Before the new one: say the Marathi for a drizzle, then the Marathi for torrential.
\end{quote}

\subsection*{You'll want to know: \mr{वळीव}}
\textbf{\mr{वळीव}} — \emph{vaḷīv} — \textquotedblleft{}a pre-monsoon storm\textquotedblright{}.

\begin{quote}
\mr{उन्हाळ्याच्या} \mr{शेवटी} \mr{वळवाचा} \mr{पाऊस} \mr{येतो}.
\end{quote}

\begin{quote}
\emph{unhāḷyācyā śevṭī vaḷvācā pāūs yeto.}
\end{quote}

\begin{quote}
At the end of summer the pre-monsoon storms come.
\end{quote}

Masculine.

Before an ending \mr{वळीव} becomes \mr{वळवा}-: \mr{वळवाचा} \mr{पाऊस}.

\begin{cousinweb}
\textbf{\mr{वळीव}} (\emph{vaḷīv}) names the sudden thunderstorm of late summer before the monsoon; its origin is not settled.
\end{cousinweb}

\begin{grammarlens}[title={What you've built}]
One more word for getting around.
\end{grammarlens}

\subsection*{Your turn}
\begin{itemize}
\raggedright
  \item \emph{Say it:} \emph{vaḷīv}
  \item \emph{Say it:} \emph{vaḷīv}, once more
  \item \emph{Say it:} \emph{unhāḷyācyā śevṭī vaḷvācā pāūs yeto.} — At the end of summer the pre-monsoon storms come.
  \item \emph{Recall:} say the Marathi for a drizzle, then the Marathi for torrential, then say \emph{vaḷīv} again
\end{itemize}

\begin{tcolorbox}[breakable,colback=teal!4,colframe=teal!35!black,title={Before you move on}]
What does \mr{वळीव} mean? (\textquotedblleft{}A pre-monsoon storm\textquotedblright{}.) Say it once more.
\end{tcolorbox}
\section[\texorpdfstring{\mr{बरसणे} (barasṇe)}{बरसणे (barasṇe)}]{\mr{बरसणे} (barasṇe) — to pour down}

\label{lesson:MR-C438-barasne}



\begin{quote}
Before the new one: say the Marathi for torrential, then the Marathi for a pre-monsoon storm.
\end{quote}

\subsection*{You'll want to know: \mr{बरसणे}}
\textbf{\mr{बरसणे}} — \emph{barasṇe} — \textquotedblleft{}to pour down\textquotedblright{}.

\begin{quote}
\mr{ढग} \mr{गडगडले} \mr{आणि} \mr{पाऊस} \mr{बरसला}.
\end{quote}

\begin{quote}
\emph{ḍhag gaḍgaḍle āṇi pāūs baraslā.}
\end{quote}

\begin{quote}
The clouds rumbled and the rain poured down.
\end{quote}

Past (he …): \mr{तो} \mr{बरसला}, \emph{to baraslā}.

Rain \mr{बरसतो}; it sounds more poetic than \mr{पडणे}.

\begin{cousinweb}
\textbf{\mr{बरसणे}} (\emph{barasṇe}) belongs with Sanskrit \emph{varṣati}, \textquotedblleft{}it rains\textquotedblright{}, where CDIAL lists Marathi \emph{varasṇẽ} (CDIAL 11394); the form with b- matches Hindi \emph{barasnā}, and how it reached Marathi is not settled.
\end{cousinweb}

\begin{grammarlens}[title={What you've built}]
One more word for getting around.
\end{grammarlens}

\subsection*{Your turn}
\begin{itemize}
\raggedright
  \item \emph{Say it:} \emph{barasṇe}
  \item \emph{Say it:} \emph{barasṇe}, once more
  \item \emph{Say it:} \emph{ḍhag gaḍgaḍle āṇi pāūs baraslā.} — The clouds rumbled and the rain poured down.
  \item \emph{Recall:} say the Marathi for torrential, then the Marathi for a pre-monsoon storm, then say \emph{barasṇe} again
\end{itemize}

\begin{tcolorbox}[breakable,colback=teal!4,colframe=teal!35!black,title={Before you move on}]
What does \mr{बरसणे} mean? (\textquotedblleft{}To pour down\textquotedblright{}.) Say it once more.
\end{tcolorbox}
\section[\texorpdfstring{\mr{भिजणे} (bhijṇe)}{भिजणे (bhijṇe)}]{\mr{भिजणे} (bhijṇe) — to get wet}

\label{lesson:MR-C438-bhijne}



\begin{quote}
Before the new one: say the Marathi for a pre-monsoon storm, then the Marathi for to pour down.
\end{quote}

\subsection*{You'll want to know: \mr{भिजणे}}
\textbf{\mr{भिजणे}} — \emph{bhijṇe} — \textquotedblleft{}to get wet\textquotedblright{}.

\begin{quote}
\mr{छत्री} \mr{नव्हती} \mr{म्हणून} \mr{मी} \mr{भिजलो}.
\end{quote}

\begin{quote}
\emph{chatrī navhatī mhaṇūn mī bhijlo.}
\end{quote}

\begin{quote}
I had no umbrella, so I got wet.
\end{quote}

Past (he …): \mr{तो} \mr{भिजला}, \emph{to bhijlā}.

A woman says \mr{मी} \mr{भिजले}; \mr{ओलाचिंब} \mr{भिजणे}, to get soaked.

\begin{cousinweb}
\textbf{\mr{भिजणे}} (\emph{bhijṇe}) is inherited: CDIAL lists Marathi \emph{bhij̈ṇẽ}, \textquotedblleft{}to be wet\textquotedblright{}, under \emph{*bhiyajyate}, \textquotedblleft{}is smeared, gets wet\textquotedblright{} (CDIAL 9502).
\end{cousinweb}

\begin{grammarlens}[title={What you've built}]
One more word for getting around.
\end{grammarlens}

\subsection*{Your turn}
\begin{itemize}
\raggedright
  \item \emph{Say it:} \emph{bhijṇe}
  \item \emph{Say it:} \emph{bhijṇe}, once more
  \item \emph{Say it:} \emph{chatrī navhatī mhaṇūn mī bhijlo.} — I had no umbrella, so I got wet.
  \item \emph{Recall:} say the Marathi for a pre-monsoon storm, then the Marathi for to pour down, then say \emph{bhijṇe} again
\end{itemize}

\begin{tcolorbox}[breakable,colback=teal!4,colframe=teal!35!black,title={Before you move on}]
What does \mr{भिजणे} mean? (\textquotedblleft{}To get wet\textquotedblright{}.) Say it once more.
\end{tcolorbox}
